% Capitolo implementativo aggiornato al 23 settembre 2026.
% Da includere nel documento principale. Nessuna compilazione eseguita.

Il presente capitolo descrive la realizzazione in MATLAB del sistema
progettato nel capitolo precedente. L'attenzione è rivolta all'organizzazione
del software, alle strutture dati e alle modifiche introdotte per correggere
le coordinate, rendere raggiungibili le destinazioni e ridurre i tempi di
calcolo. I listing riportano estratti delle operazioni principali; non
costituiscono programmi autonomi quando omettono i controlli circostanti.

Il flusso corrente assegna gli ordini prima dell'esecuzione e mantiene
un unico stato temporale. Le finestre organizzano l'osservazione di una
dinamica per eventi: non generano nuove assegnazioni e non ricostruiscono
la rete a ogni chiamata.

\section{Organizzazione del software}
\label{sec:impl-organizzazione}

Il punto di ingresso è \texttt{main.m}. La cartella \texttt{importExcel}
contiene la lettura dei dati; \texttt{utility} comprende preparazione,
ordinamento, storici e grafici. Le funzioni di compatibilità, assegnazione
e generazione delle matrici sono in \texttt{modelGeneration}; quelle
spaziali in \texttt{pathPlan}. La cartella \texttt{temporal} contiene
inizializzazione, avanzamento e riepilogo, mentre \texttt{tests} raccoglie
le verifiche.

Nel repository rimangono anche moduli della precedente esecuzione per batch.
La loro presenza non implica che vengano richiamati dal programma principale:
il flusso attuale utilizza il gestore temporale e le matrici generate per
l'insieme delle missioni ammesse.

\subsection{Configurazione e ciclo principale}

\begin{lstlisting}[language=Matlab,caption={Configurazione corrente del simulatore.},label={lst:impl-parametri}]
numberOfResources = 8;
missionTimeLimit = 180;
windowDuration = 15*60;
maxT = 8*3600;
showPlots = true;
\end{lstlisting}

Lo script importa \texttt{1su11.xlsx}, costruisce la mappa dalle immagini,
prepara le missioni e assegna gli ordini. Inizializza quindi il pianificatore
A* e lo stato temporale. I percorsi dei file sono costruiti rispetto alla
cartella dello script.

\begin{lstlisting}[language=Matlab,caption={Avanzamento sullo stesso stato persistente.},label={lst:impl-ciclo}]
tic;
while any(ismember(simulationState.Missions.Status, ...
        ["pending","running"]))
    simulationState = advanceTemporalSimulation( ...
        simulationState, ...
        simulationState.Clock + windowDuration, ...
        planner, occupancyMatrix);
end
totalElapsedTime = toc;
\end{lstlisting}

\texttt{totalElapsedTime} misura il ciclo sul computer, escludendo la
preparazione precedente e i grafici finali. Il tempo simulato è invece
conservato in \texttt{Clock}. \texttt{maxT} è utilizzato nel riepilogo,
non nella condizione di prosecuzione del ciclo.

\section{Importazione e preparazione delle missioni}
\label{sec:impl-importazione}

\texttt{importFromExcel} normalizza le intestazioni eliminando gli spazi
esterni e convertendole in maiuscolo. Per l'ordine cerca, in sequenza,
\texttt{ID\_ORDER}, \texttt{OID\_ORDER} e \texttt{CD\_ORDER}, usando la
prima colonna disponibile. Il risultato viene esposto internamente come
\texttt{ID\_ORDER}; non vengono combinati valori di colonne differenti.
Questa gestione conserva l'informazione anche quando l'esportazione usa
\texttt{OID\_ORDER}.

Gli identificativi vengono convertiti in stringhe e i valori mancanti
normalizzati a stringhe vuote. In assenza di una colonna riconosciuta viene
emesso un avviso e le missioni sono trattate come indipendenti. La priorità
proviene da \texttt{NM\_PRIORITY}, con valore predefinito 1 per gli elementi
mancanti o per l'intera colonna assente.

\subsection{Identità e ordinamento stabile}

\texttt{defineMissions} costruisce la tabella con codice, priorità,
coordinate FROM e TO, ordine e \texttt{missionUid}. L'UID identifica la
riga di origine, resta invariato dopo ordinamento e filtraggio e deve essere
univoco. Il codice commerciale può invece ripetersi senza accorpare attività
distinte.

Le coordinate sono contenute nelle sottotabelle \texttt{missionFrom}
e \texttt{missionTo}. \texttt{groupMissionsByOrder} raggruppa gli ordini
secondo la prima comparsa, assegna gruppi distinti alle missioni senza ordine
e ordina mediante gruppo, priorità negativa e posizione originale. Si
ottiene così una priorità decrescente interna all'ordine con parità stabili.

\subsection{Correzione delle coordinate esterne}

Prima della validazione resta attiva l'euristica che divide per 1000 e
arrotonda le coordinate superiori a 999. Questa regola è legata alla
rappresentazione dei dati disponibili e non è una conversione universale
delle unità. Successivamente i valori finiti esterni vengono proiettati
entro i limiti della mappa.

\begin{lstlisting}[language=Matlab,caption={Estratto della proiezione delle coordinate.},label={lst:impl-coordinate}]
limits = [maxX maxY maxX maxY];
for k = 1:numel(coordinateNames)
    values = workFlow.(coordinateNames{k});
    outside = isfinite(values) & ...
        (values < 1 | values > limits(k));
    values(outside) = min(max(values(outside),1),limits(k));
    workFlow.(coordinateNames{k}) = values;
end
\end{lstlisting}

La modifica interessa le coordinate operative in memoria, senza riscrivere
l'Excel. Un contatore segnala le missioni corrette. Le coordinate non finite
rimangono invalide e le righe corrispondenti sono conservate in
\texttt{discardedMissions}.

Sul file corrente sono state verificate 1105 righe: la validazione precedente
ne escludeva 120 per superamento dei limiti; la proiezione recupera quelle
120 e ammette tutte le 1105 nella preparazione. Questa verifica non equivale
al completamento di tutte le missioni.

\section{Assegnazione iniziale e matrici della rete}
\label{sec:impl-assegnazione}

\texttt{defineAttitude} definisce i limiti delle risorse e
\texttt{assignMission} ne valuta la compatibilità con ciascuna missione.
\texttt{assignOrders} interseca gli insiemi compatibili per tutte le
missioni dell'ordine e sceglie il proprietario con il minor numero di
missioni già assegnate.

\begin{lstlisting}[language=Matlab,caption={Estratto della scelta della risorsa proprietaria.},label={lst:impl-owner}]
candidates = 1:numberOfResources;
for k = rows'
    candidates = intersect(candidates,eligible{k});
end
% Prima della scelta viene controllata l'assenza di candidati.
[~,leastLoaded] = min(load(candidates));
owner = candidates(leastLoaded);
missions.assignedResource(rows) = owner;
load(owner) = load(owner) + numel(rows);
\end{lstlisting}

Un insieme vuoto di candidati produce un errore. A parità di carico prevale
l'indice minore. La tabella restituita registra ordine, risorsa e cardinalità;
la scelta non viene ricalcolata durante l'avanzamento sulla base dei tempi
accumulati.

\texttt{generateMatrices} costruisce tre transizioni e due posti interni
per missione, oltre al posto comune delle risorse. Per le missioni successive
dello stesso ordine aggiunge posti di collegamento e rimuove il rilascio
e il nuovo prelievo intermedi dal posto comune. La risorsa mantiene quindi
il proprio colore lungo la sequenza.

Le matrici sono generate inizialmente in forma densa e convertite in forma
sparsa nell'inizializzazione temporale. La marcatura ha una colonna per
risorsa; all'inizio ogni colore possiede un token nel posto comune.
L'abilitazione confronta la marcatura del colore con gli ingressi della
transizione, mentre lo scatto sottrae gli ingressi e aggiunge le uscite.
L'ottimizzazione spaziale descritta in seguito non modifica il costo di
costruzione delle matrici della rete.
\section{Stato persistente e avanzamento temporale}
\label{sec:impl-stato}

\texttt{initializeTemporalSimulation} crea la struttura con missioni,
ordini, matrici, marcatura, percorsi e storici. Nella tabella delle missioni
sono aggiunti \texttt{Status}, \texttt{StartTime},
\texttt{ScheduledFinish}, \texttt{FinishTime}, \texttt{Duration},
\texttt{RemainingTime} e \texttt{FailureReason}. Le missioni ammesse
iniziano in stato \texttt{pending}, quelle escluse in \texttt{excluded}.
Le escluse rimangono nel rendiconto ma non creano sottoreti operative.

La tabella degli ordini mantiene proprietario, conteggi, budget, deadline,
completamento e ritardo. Le chiavi logiche distinguono anche le missioni
senza ordine. \texttt{Plans} conserva i percorsi; \texttt{ResourcePos}
le posizioni delle risorse, inizialmente non definite. \texttt{Events}
registra tempo, UID, risorsa, transizione e tipo di evento.
\texttt{Windows} raccoglie gli estremi delle finestre, le missioni avviate,
completate, fallite e ancora in esecuzione, con i relativi residui.

\subsection{Avvii e completamenti}

\texttt{advanceTemporalSimulation} completa anzitutto le missioni il cui
termine programmato è maturato all'orologio corrente. Aggiorna marcatura,
storico e posizione delle risorse. Se non è stato raggiunto il limite della
finestra, cerca per ciascuna risorsa la prima missione pendente abilitata e
ne tenta la pianificazione. Un percorso valido determina l'avvio:

\begin{lstlisting}[language=Matlab,caption={Registrazione dei tempi all'avvio della missione.},label={lst:impl-avvio}]
state.Missions.Status(row) = "running";
state.Missions.StartTime(row) = state.Clock;
state.Missions.Duration(row) = planResult.Time;
state.Missions.ScheduledFinish(row) = ...
    state.Clock + planResult.Time;
\end{lstlisting}

L'orologio passa poi al minimo tra il prossimo completamento e il limite
della finestra. Non viene effettuata un'attesa proporzionale ai secondi
simulati. Un completamento esattamente al confine viene registrato; la
chiamata successiva riprende dallo stesso istante, senza aggiungere un
ritardo artificiale al nuovo avvio.

Se non restano missioni in corso ma sono presenti missioni assegnate non
abilitabili, viene generata un'asserzione di deadlock. Lo stato
\texttt{blocked} è previsto nei riepiloghi, ma non viene attribuito
automaticamente da questa gestione. Al termine dell'avanzamento vengono
aggiornati ordini, residui e finestra e viene verificata la conservazione
di un token per risorsa.

\subsection{Budget e deadline}

Il budget ordinario è il numero di missioni dell'ordine moltiplicato per
\texttt{missionTimeLimit}, includendo anche le righe escluse conservate
nel rendiconto. La tabella \texttt{orderDurationOverrides}, con colonne
\texttt{OrderId} e \texttt{MaxDuration}, può sostituire il budget di un
ordine. Valori non positivi o non finiti e override duplicati sono respinti.
La deadline è fissata soltanto al primo avvio:

\begin{lstlisting}[language=Matlab,caption={Definizione della deadline all'avvio effettivo.},label={lst:impl-deadline}]
if isnan(state.Orders.StartTime(o))
    state.Orders.StartTime(o) = state.Clock;
    state.Orders.Deadline(o) = ...
        state.Clock + state.Orders.MaxDuration(o);
end
\end{lstlisting}

\texttt{refreshTemporalOrders} considera completo un ordine soltanto
quando tutte le missioni sono completate. La presenza di missioni fallite,
escluse o bloccate lo rende incompleto; negli altri casi distingue ordine
aperto e non avviato. Il ritardo usa il completamento finale oppure, per
ordini non conclusi con deadline definita, l'orologio corrente.

\texttt{summarizeTemporalSimulation} conta missioni e ordini per stato,
ordini puntuali e in ritardo e missioni oltre il riferimento nominale.
L'indicatore di completamento entro il turno richiede sia tutti gli ordini
completati sia l'orologio entro \texttt{maxT}. Questi controlli valutano
il risultato e non troncano le esecuzioni.

\section{Mappa, coordinate e pianificazione}
\label{sec:impl-percorsi}

\texttt{computeMap} legge le immagini del layout. Per immagini a colori
seleziona il canale verde; i pixel diversi da 255 diventano ostacoli.
La mappa binaria corrente ha dimensioni nominali $50.2\times95.1$ e
risoluzione 10, producendo una griglia $951\times502$. La formula temporale
non utilizza direttamente queste dimensioni fisiche, ma il numero di punti
del percorso.

\subsection{Indici ai bordi della griglia}

\texttt{scalePoints} arrotonda le coordinate e inverte Y. La formula
precedente produceva indice zero quando Y coincideva con il massimo;
la correzione è l'aggiunta dell'unità richiesta dagli indici MATLAB:

\begin{lstlisting}[language=Matlab,caption={Conversione corretta delle coordinate operative.},label={lst:impl-indici}]
startPoint = [round(startP(1)), ...
    round(mapSize(2) + 1 - startP(2))];
endPoint = [round(endP(1)), ...
    round(mapSize(2) + 1 - endP(2))];
\end{lstlisting}

Y uguale a 502 diventa così colonna 1. Il problema era emerso anche nella
missione 6004972 dopo la proiezione preliminare sul bordo.
\texttt{pathPlanning} verifica comunque che posizione della risorsa,
FROM e TO siano finiti, interi e interni alla mappa.

\subsection{Estremi raggiungibili e continuità}

Il pianificatore legge la propria matrice di occupazione e ne confronta le
dimensioni con quelle fornite. La connettività dipende dalla proprietà
\texttt{DiagonalSearch}. Una risorsa senza posizione precedente viene
collocata al FROM iniziale; una posizione occupata è riportata alla cella
libera più vicina. Un FROM occupato è corretto nella componente raggiungibile
dalla risorsa. Se il FROM è libero ma isolato, l'avvicinamento può fallire.

La prima chiamata ad A* calcola il percorso dalla risorsa al FROM; il suo
ultimo punto diventa il FROM effettivo. Il TO è poi ricercato nella stessa
componente anche quando la cella richiesta è libera:

\begin{lstlisting}[language=Matlab,caption={Estratto della correzione del TO e della seconda tratta.},label={lst:impl-to}]
fromPoint = path(end,:);
toPoint = findNearestZero( ...
    plannerMatrix,toPoint,fromPoint,connectivity);
% Il codice verifica che toPoint non sia vuoto.
secondPath = plan(planner,fromPoint,toPoint);
\end{lstlisting}

La precedente correzione condizionata alla sola occupazione non gestiva i
TO liberi ma isolati. Nella prova della missione 6003316, UID 1037, il TO
$[905,502]$ è stato sostituito con $[905,499]$, espresso come riga e colonna.
Il percorso è diventato realizzabile e la durata della prova è risultata
208 s. La risorsa era inizialmente presso il FROM; in una sequenza completa
un diverso avvicinamento può modificare tale durata.

\subsection{Formula della durata}

Le due tratte vengono concatenate e il tempo viene stimato dal numero
di punti risultante:

\begin{lstlisting}[language=Matlab,caption={Durata simulata del percorso concatenato.},label={lst:impl-durata}]
path = [path;secondPath];
timePath = size(path,1)/5.5;
time = round(timePath + timePath*0.2) + 60;
\end{lstlisting}

Il conteggio include il punto comune nella concatenazione e non distingue
la lunghezza dei passi diagonali. I coefficienti e il contributo fisso di
60 s sono parte del modello empirico e richiedono calibrazione per una
corrispondenza fisica. Non vengono modellate collisioni spazio-temporali
tra i mezzi: la mappa usata da A* contiene ostacoli statici.

\section{Fallimenti e aggiornamento delle risorse}
\label{sec:impl-fallimenti}

Un errore sugli estremi, l'assenza di celle utilizzabili o una tratta
irrealizzabile producono una tabella di pianificazione vuota. Il gestore
assegna \texttt{failed} e registra l'evento. Per proseguire, sottrae dalla
marcatura gli ingressi della transizione iniziale e aggiunge le uscite
della transizione finale del ramo fallito. Il token resta riservato se
seguono missioni dell'ordine e torna al posto comune al termine della sequenza.
Non vengono aggiornati posizione o lavoro accumulato per una missione fallita.

\texttt{updateBatch} interviene al completamento riuscito: aggiorna tempo,
codici missione, ordini e priorità nello storico e pone la risorsa
nell'ultima cella del percorso. Il gestore aggiunge anche gli UID, preservando
l'identità delle occorrenze con codici uguali.

\texttt{selectedBatch} espone lo storico per risorsa, non una nuova
configurazione scelta mediante funzione di costo. Le tabelle finali delle
missioni fallite e bloccate sono estrazioni della cronologia temporale.
\section{Diagnostica grafica dei percorsi falliti}
\label{sec:impl-diagnostica}

La firma di \texttt{pathPlanning} comprende un secondo output diagnostico
con i campi \texttt{Reason}, \texttt{Start}, \texttt{From}, \texttt{To},
\texttt{RequestedFrom} e \texttt{RequestedTo}. Il gestore lo conserva
nel campo \texttt{PathDiagnostics}, associandolo alla riga della missione.
Il motivo generale nella tabella rimane \texttt{no\_feasible\_path};
la struttura aggiuntiva distingue invece la causa specifica.

La diagnostica viene prodotta anche per una pianificazione riuscita.
In questo contesto il motivo \texttt{completed} indica il successo del
calcolo del percorso, non il completamento temporale della missione, che
passa prima a \texttt{running}. I punti richiesti sono già indici della
griglia e possono aver subito la proiezione iniziale: non sono una copia
delle coordinate grezze dell'Excel.

\texttt{plotFailedMissionPaths} filtra le missioni fallite e apre una
finestra con menu di selezione per codice, UID e ordine. Mostra gli ostacoli,
il FROM verde, il TO rosso e la posizione della risorsa blu, distinguendo
estremi richiesti ed effettivi. L'asse orizzontale rappresenta la colonna,
quello verticale la riga. I segmenti tratteggiati sono collegamenti
indicativi e non percorsi validati.

In assenza della diagnostica dettagliata vengono mostrati soltanto i punti
richiesti, segnalando la limitazione. Se non esistono missioni fallite,
la funzione non crea una figura. Il richiamo è integrato nei grafici finali
del programma principale.

\section{Statistiche delle durate e fasce temporali}
\label{sec:impl-statistiche}

Il file separato \texttt{calcolaTempiMissioni.m} analizza una simulazione
già disponibile, selezionando le sole missioni completate con durata
finita e non negativa:

\begin{lstlisting}[language=Matlab,caption={Filtro delle durate utilizzate nell'analisi.},label={lst:impl-filtro}]
valide = string(missionTimeline.Status)=="completed" & ...
    isfinite(missionTimeline.Duration) & ...
    missionTimeline.Duration>=0;
durate = missionTimeline.Duration(valide);
\end{lstlisting}

La funzione restituisce conteggi, minimo, media, mediana e massimo, in
secondi e minuti. La colonna \texttt{MissioniEscluse} conta tutte le righe
non usate nell'analisi; non coincide necessariamente con lo scarto preliminare
in \texttt{discardedMissions}. I codici duplicati non vengono accorpati.

Il secondo output \texttt{fasce} riporta estremi e conteggi degli intervalli.
L'ampiezza predefinita è 30 s ed è personalizzabile. Il bordo superiore
complessivo è posto oltre il massimo osservato; una durata su un confine
interno appartiene alla fascia successiva. Il grafico a barre comprende
etichette, conteggi e fasce intermedie vuote.

\begin{lstlisting}[language=Matlab,caption={Analisi con fasce di trenta secondi.},label={lst:impl-statistiche}]
[statistiche, fasce, figura] = ...
    calcolaTempiMissioni(missionTimeline,30);
\end{lstlisting}

Senza durate valide le statistiche valgono \texttt{NaN}, la tabella delle
fasce è vuota e non viene creata una figura. L'analisi può guidare la scelta
del budget, ma non modifica automaticamente le deadline.

\section{Ottimizzazione delle componenti raggiungibili}
\label{sec:impl-cache}

\subsection{Costo della ricerca precedente}

La versione precedente di \texttt{findNearestZero} ricostruiva la componente
raggiungibile con una visita in ampiezza a ogni chiamata vincolata a un
riferimento. Dopo l'estensione del controllo del TO a tutte le missioni,
la stessa regione della mappa era esplorata ripetutamente. L'ottimizzazione
ha riguardato questa operazione, distinta dal calcolo dei percorsi A*.

\subsection{Dati persistenti e invalidazione}

La funzione conserva la maschera delle celle libere, la connettività,
le etichette delle componenti e le coordinate dei loro punti. La cache
viene invalidata quando cambia la connettività o la maschera:

\begin{lstlisting}[language=Matlab,caption={Estratto dell'invalidazione della cache spaziale.},label={lst:impl-cache}]
persistent cachedFree cachedConnectivity ...
    componentLabels componentPoints

if isempty(cachedConnectivity) || ...
        cachedConnectivity ~= connectivity || ...
        ~isequal(cachedFree,free)
    cachedFree = free;
    cachedConnectivity = connectivity;
    componentLabels = zeros(mapSize,'uint32');
    componentPoints = cell(0,1);
end
\end{lstlisting}

Viene conservata soltanto la mappa più recente. Il confronto rileva anche
cambiamenti degli ostacoli a dimensioni invariate. Variazioni numeriche in
celle già occupate non richiedono una ricostruzione se rimane identica
la classificazione libero/occupato.

Se la cella di riferimento non è etichettata, la visita in ampiezza esplora
la sua componente e ne memorizza i punti. Le chiamate successive riutilizzano
i dati. Un target già nella stessa componente è restituito direttamente;
negli altri casi si calcola la distanza quadratica dai punti memorizzati.
Le parità vengono risolte per riga e colonna, senza dipendere dall'ordine
di visita della coda.

La ricerca senza riferimento considera tutte le celle libere. Un riferimento
occupato produce un risultato vuoto. Il confronto della maschera e la
ricerca del minimo mantengono un costo: la cache elimina l'esplorazione
ripetuta della componente, non ogni scansione. Non vengono memorizzati i
percorsi A*.

Per azzerare esplicitamente la cache durante una verifica si può usare
\texttt{clear findNearestZero}. L'azzeramento delle sole variabili del
programma principale non cancella le variabili persistenti della funzione.

\section{Verifiche funzionali e prestazioni}
\label{sec:impl-verifiche}

Le verifiche sono raccolte in funzioni MATLAB nella cartella \texttt{tests}.
Le regressioni collegate alle modifiche comprendono:
\begin{itemize}
    \item \texttt{testMissionPreparation}: proiezione, identità delle
          missioni ed esclusione dei valori non finiti;
    \item \texttt{testScalePoints}: conversione dei bordi e percorso con
          estremi sui limiti della mappa;
    \item \texttt{testNearestReachableZero}: distanza minima, parità,
          connettività e TO libero ma isolato;
    \item \texttt{testNoFeasiblePath}: fallimenti senza completamenti
          fittizi e senza aggiornamenti impropri della posizione;
    \item \texttt{testFailedMissionPaths}: estremi diagnostici, selezione
          tra codici duplicati e assenza di fallimenti;
    \item \texttt{testNearestZeroCache}: invalidazione ed equivalenza
          rispetto a una ricerca indipendente.
\end{itemize}

Il test della cache modifica gli ostacoli a dimensione invariata, alterna
connettività e interroga componenti differenti. Esegue anche 320 confronti
su mappe generate con seme riproducibile. Il riferimento ricostruisce la
raggiungibilità con espansione iterativa della maschera e applica gli stessi
criteri di distanza e parità.

Il progetto comprende inoltre test sullo stato temporale e sulla riserva
delle risorse per ordine. La loro presenza va distinta dalla riesecuzione
sistematica dopo ogni modifica: per la cache sono stati eseguiti il nuovo
test e le regressioni pertinenti su celle, percorsi falliti, diagnostica
e conversione dei bordi.

\subsection{Confronto prima e dopo l'ottimizzazione}

Il confronto ha utilizzato le prime venti missioni preparate dal file
corrente, lo stesso pianificatore e una posizione iniziale non definita
per ogni chiamata. Sono stati confrontati integralmente percorso e
diagnostica: punti, durate e risultati sono rimasti identici. La cache
della nuova funzione era vuota all'inizio della misura.

\begin{table}[htbp]
\centering
\begin{tabular}{lr}
\toprule
Versione, senza profiler & Tempo per 20 missioni\\
\midrule
Ricerca precedente & 4.371 s\\
Ricerca con cache inizialmente vuota & 1.674 s\\
\bottomrule
\end{tabular}
\caption{Confronto sul campione usato per verificare la cache.}
\label{tab:impl-benchmark}
\end{table}

Il rapporto dei tempi è circa $2.61$. La misura riguarda soltanto il campione
di pianificazione e non l'intera simulazione con risorse aggiornate lungo
la sequenza. Non è quindi un fattore garantito per il programma completo.

Una misura precedente con profiler attribuiva alla ricerca delle celle
circa 72.7 s su 74.4 s totali sul campione. Il profiler introduce overhead
rilevante nei cicli con molte chiamate: tali valori non sono direttamente
confrontabili con la Tabella~\ref{tab:impl-benchmark}. Il risultato utile
è la riduzione del lavoro ripetuto senza variazioni degli esiti.

\section{Risultati disponibili e uso del sistema}
\label{sec:impl-risultati}

Al termine dello script, \texttt{missionTimeline} e \texttt{orderTimeline}
espongono le cronologie delle attività e degli ordini.
\texttt{transitionEvents} raccoglie gli eventi e \texttt{executiveSummary}
il riepilogo. \texttt{miniBatchHistory} corrisponde alla tabella delle
finestre: il nome non indica un algoritmo di costruzione di batch con
cardinalità minima e massima.

Con \texttt{showPlots} attivo vengono mostrati sequenza per risorsa,
tempi e deadline e diagnostica dei fallimenti. La distribuzione delle
durate si ottiene con il richiamo separato alla funzione delle statistiche.

\begin{lstlisting}[language=Matlab,caption={Sequenza di utilizzo dopo la configurazione.},label={lst:impl-utilizzo}]
main
[statistiche, fasce, figura] = ...
    calcolaTempiMissioni(missionTimeline,30);

% Per riaprire la diagnostica:
plotFailedMissionPaths(simulationState,occupancyMatrix);
\end{lstlisting}

Per cambiare il budget nominale si modifica \texttt{missionTimeLimit}
e si riesegue la simulazione; gli override restano specifici per ordine.
La durata della finestra controlla la frequenza delle osservazioni e non
introduce una riassegnazione del lavoro.

Le strutture dati mantengono distinte preparazione, fattibilità geometrica,
esecuzione e valutazione del servizio. Una riga recuperata dalla proiezione
non è automaticamente una missione completata; una pianificazione riuscita
non è ancora un evento di completamento; l'esaurimento del lavoro operativo
può lasciare un ordine incompleto. Tale separazione permette di interpretare
correttamente gli indicatori e i grafici prodotti.